\section{Objectifs}

\begin{itemize}
\item Reconnaître et exploiter les fonctions affines et les fonctions de référence du programme.
\item Relier expression algébrique, courbe, signe et variations.
\item Construire et exploiter un tableau de signes d'un produit ou d'un quotient.
\item Relier sens de variation, coefficient directeur et droite représentative d'une fonction affine.
\item Résoudre graphiquement, algébriquement ou numériquement des équations et inéquations liées aux fonctions de référence.
\item Déterminer ou approcher un extrémum et traiter un problème d'optimisation.
\end{itemize}

\section{Prérequis}

\begin{itemize}
\item fonctions : vocabulaire, images et antécédents
\item équations et inéquations
\item repère et droites
\end{itemize}

\section{Activité de découverte}

On compare plusieurs programmes de calcul. Leurs courbes n'ont pas la même forme : droite, parabole, hyperbole, courbe en V, racine carrée ou cube. Ces formes deviennent des repères visuels, mais la résolution d'un problème doit aussi pouvoir s'appuyer sur l'algèbre et les variations.

\section{Vocabulaire}

\begin{itemize}
\item \textbf{croissante}, \textbf{décroissante}, \textbf{monotone}
\item \textbf{maximum}, \textbf{minimum}, \textbf{extrémum}
\item \textbf{tableau de variations}, \textbf{tableau de signes}
\item \textbf{fonction affine}, \textbf{coefficient directeur}, \textbf{ordonnée à l'origine}, \textbf{taux d'accroissement}
\item \textbf{fonction valeur absolue}, \textbf{carré}, \textbf{inverse}, \textbf{racine carrée}, \textbf{cube}
\end{itemize}

\section{Cours}

\subsection{Fonctions affines}

Une fonction affine s'écrit \texttt{f(x)=mx+p}.

\begin{itemize}
\item \texttt{m} est le coefficient directeur et le taux d'accroissement : pour \texttt{x1≠x2}, \texttt{(f(x2)-f(x1))/(x2-x1)=m} ;
\item \texttt{p=f(0)} est l'ordonnée à l'origine ;
\item si \texttt{m>0}, la fonction est croissante ;
\item si \texttt{m<0}, elle est décroissante ;
\item si \texttt{m=0}, elle est constante.
\end{itemize}

Le signe d'une fonction affine se lit à partir de son éventuelle racine et du signe de \texttt{m}.

\subsection{Fonctions de référence}

\begin{itemize}
\item \texttt{x ↦ |x|} est définie sur \texttt{ℝ}, positive, décroissante sur \texttt{]-∞;0]} puis croissante sur \texttt{[0;+∞[} ; son minimum vaut \texttt{0}.
\item \texttt{x ↦ x²} est définie sur \texttt{ℝ}, positive, décroissante sur \texttt{]-∞;0]} puis croissante sur \texttt{[0;+∞[} ; son minimum vaut \texttt{0}.
\item \texttt{x ↦ 1/x} est définie sur \texttt{ℝ\{0}} et décroissante sur chacun des intervalles \texttt{]-∞;0[} et \texttt{]0;+∞[}.
\item \texttt{x ↦ √x} est définie sur \texttt{[0;+∞[} et croissante.
\item \texttt{x ↦ x³} est définie sur \texttt{ℝ} et croissante.
\end{itemize}

\subsection{Signe d'un produit ou d'un quotient}

Pour un produit ou un quotient, on détermine d'abord les zéros des facteurs et les valeurs interdites du dénominateur, puis on combine les signes dans un tableau.

Le tableau de signes permet de résoudre \texttt{f(x)>0}, \texttt{f(x)≤0} lorsque \texttt{f} est un produit ou un quotient d'expressions simples.

\subsection{Variations et extrémums}

Une fonction est croissante sur un intervalle lorsque l'ordre des nombres de départ est conservé par leurs images. Elle est décroissante lorsque cet ordre est inversé.

Un tableau de variations synthétise ces informations. Un maximum ou un minimum doit toujours être indiqué avec l'ensemble ou l'intervalle sur lequel il est considéré.

\subsection{Résoudre et comparer}

Selon la situation, on choisit le registre le plus efficace :

\begin{itemize}
\item graphique pour lire ou estimer les solutions de \texttt{f(x)=k}, \texttt{f(x)<k}, \texttt{f(x)=g(x)} ;
\item algébrique lorsque les expressions s'y prêtent ;
\item calculatrice, logiciel ou Python pour obtenir une approximation ou explorer un problème.
\end{itemize}

Pour les fonctions affine, valeur absolue, carré, inverse, racine carrée et cube, les équations ou inéquations simples avec une constante doivent pouvoir être traitées graphiquement ou algébriquement.

\subsection{Optimisation}

Un problème d'optimisation consiste à chercher la plus grande ou la plus petite valeur d'une grandeur modélisée par une fonction. La démarche comprend : choisir la variable, déterminer son domaine, construire la fonction, étudier ou explorer ses variations, puis interpréter l'extrémum dans le contexte.

\section{Exemples corrigés}

\begin{itemize}
\item Pour \texttt{f(x)=2x-6}, \texttt{f(x)>0} équivaut à \texttt{x>3}.
\item \texttt{x²=4} a deux solutions : \texttt{-2} et \texttt{2}.
\item \texttt{√x<3} avec \texttt{x≥0} équivaut à \texttt{0≤x<9}.
\item Pour \texttt{f(x)=-3x+2}, le coefficient directeur est négatif : \texttt{f} est décroissante sur \texttt{ℝ}.
\end{itemize}

\coursefigure{/images/mathematiques/latex-migration/2nde-fonctions-reference-variations-figure-1.svg}{Repère montrant une droite, une parabole et une courbe inverse stylisée.}{Les courbes de référence sont des repères ; la résolution doit aussi pouvoir changer de registre.}

\section{Algorithmique en Python}

\subsection{Approcher un minimum par balayage}

\begin{verbatim}
def f(x):
    return (x - 2)**2 + 1

meilleur_x = -5
for k in range(-50, 51):
    x = k / 10
    if f(x) < f(meilleur_x):
        meilleur_x = x
print(meilleur_x, f(meilleur_x))
\end{verbatim}

Le balayage fournit une approximation. Une dichotomie ou un pas plus fin peut ensuite améliorer la précision selon le problème.

\section{Démonstrations à connaître ou à travailler}

\begin{itemize}
\item variations des fonctions affines ;
\item variations des fonctions carré et inverse ;
\item position relative de \texttt{y=x} et \texttt{y=x²} pour \texttt{x≥0}.
\end{itemize}

\section{Erreurs fréquentes}

\begin{itemize}
\item Dire que la fonction carré est croissante sur tout \texttt{ℝ}.
\item Mélanger signe de \texttt{f(x)} et sens de variation de \texttt{f}.
\item Oublier que la fonction inverse n'est pas définie en \texttt{0}.
\item Oublier que la racine carrée n'est définie que pour \texttt{x≥0}.
\item Donner un maximum ou un minimum sans préciser l'intervalle étudié.
\end{itemize}

\section{Formules et idées à retenir}

\begin{itemize}
\item \texttt{f(x)=mx+p} : le signe de \texttt{m} détermine le sens de variation.
\item \texttt{x²≥0} et \texttt{|x|≥0} pour tout réel \texttt{x}.
\item \texttt{√x} nécessite \texttt{x≥0} ; \texttt{1/x} nécessite \texttt{x≠0}.
\item Un tableau de signes et un tableau de variations répondent à deux questions différentes.
\end{itemize}

\section{Synthèse}

Le programme de seconde demande de savoir circuler entre expression, tableau, graphique et outil numérique. Les fonctions de référence fournissent des modèles connus ; les variations, les signes et les extrémums permettent ensuite de résoudre des problèmes et d'optimiser une grandeur.
